\documentclass{article}
\usepackage[margin=1in]{geometry}
\usepackage{booktabs}
\usepackage{graphicx}
\usepackage{hyperref}
\usepackage{amsmath}

\title{Core ML: Modular Long-Context Sequence Modeling}
\author{Ansh}
\date{SAiDL Summer Assignment 2026}

\begin{document}
\maketitle

\section{Overview}

This report studies long-context sequence modeling on WikiText-2 using a modular
Transformer language model. The codebase allows independent swaps of attention
mechanisms, positional encodings, and convolution-attention hybrid blocks.

\section{Baseline}

\subsection{Architecture}

The baseline is a pre-normalized causal Transformer language model with tied
token embedding and language-model head weights. The default configuration uses
a context length of 1024 tokens, embedding dimension 512, 8 heads, 6 layers, and
a feed-forward dimension of 2048.

\subsection{Training Setup}

\begin{itemize}
    \item Dataset: WikiText-2 raw.
    \item Tokenizer: GPT-2 tokenizer.
    \item Optimizer: AdamW.
    \item Model width/depth: 512 hidden dimension, 8 attention heads, 6 Transformer layers, and feed-forward dimension 2048.
    \item Sequence length: 1024 tokens for the initial baseline and attention-variant experiments.
    \item Batch size: 4.
    \item Logged metrics: training loss, validation loss, validation perplexity,
    peak GPU memory, throughput, and elapsed time.
\end{itemize}

\begin{table}[h]
\centering
\begin{tabular}{lrrrrr}
\toprule
Model & Context & Val Loss & Val PPL & Val Tokens/sec & Train Peak Mem (MB) \\
\midrule
Standard + Sinusoidal & 1024 & 5.862 & 351.41 & 29636.66 & 4182.75 \\
\bottomrule
\end{tabular}
\caption{Baseline results.}
\end{table}

The standard Transformer with sinusoidal positional encoding serves as the
reference point for all later comparisons. At step 8500 it reaches a validation
loss of 5.862 and validation perplexity of 351.41 while maintaining high
throughput and moderate memory use. This baseline is strong enough to act as a
meaningful control but still leaves ample room for architectural improvement.

\section{Attention Variants}

The implemented attention variants are standard causal self-attention, sliding
window attention, linear kernelized attention, multi-query attention, and
grouped-query attention.

\begin{table}[h]
\centering
\begin{tabular}{lrrrrrr}
\toprule
Attention & Context & Step & Val Loss & Val PPL & Val Tokens/sec & Train Peak Mem (MB) \\
\midrule
Standard & 1024 & 8500 & 5.862 & 351.41 & 29636.66 & 4182.75 \\
Sliding Window & 1024 & 5500 & 6.052 & 424.96 & 28301.10 & 4188.65 \\
Linear & 1024 & 8500 & 5.856 & 349.49 & 23573.86 & 4981.47 \\
GQA & 1024 & 8500 & 5.810 & 333.78 & 29830.34 & 4155.67 \\
\bottomrule
\end{tabular}
\caption{Attention ablation results.}
\end{table}

Among the attention variants evaluated so far, grouped-query attention (GQA)
performs best. It achieves the lowest validation loss and validation perplexity
while also remaining efficient in both throughput and memory. Linear attention
slightly improves over the standard baseline in validation perplexity, but this
comes with noticeably higher memory consumption and lower throughput. The
sliding-window model performs worst in the current results, although that run
was logged at a smaller training budget than the others and should therefore be
treated as a preliminary comparison rather than a final conclusion.

\section{Positional Encodings}

The implemented positional choices are sinusoidal positional encodings, rotary
positional embeddings, ALiBi, and learned relative position bias. In addition
to the main 1024-token context experiments, I evaluate length extrapolation by
training the best positional configuration at $L_{\mathrm{train}}=512$ and
evaluating it at $L_{\mathrm{test}}\in\{512,1024,2048\}$.

\begin{table}[h]
\centering
\begin{tabular}{lrrrrr}
\toprule
Position Method & Attention & Context & Step & Val PPL & Notes \\
\midrule
Sinusoidal & GQA & 1024 & 8500 & 333.78 & Fully completed run \\
RoPE & GQA & 1024 & 6500 & 150.28 & Best logged validation checkpoint \\
RoPE (final) & GQA & 1024 & 8500 & 152.73 & Completed rerun, slight late-stage degradation \\
ALiBi & GQA & 1024 & 7000 & 146.24 & Best logged validation checkpoint \\
ALiBi (final) & GQA & 1024 & 8500 & 147.87 & Completed run, slight late-stage degradation \\
Relative Bias & GQA & 1024 & 8500 & 235.86 & Completed run, weaker than ALiBi/RoPE \\
\bottomrule
\end{tabular}
\caption{Current positional encoding results.}
\end{table}

ALiBi is the strongest positional encoding tested so far in the current setup.
The completed GQA + ALiBi run achieved its best logged validation performance at
step 7000, reaching validation loss 4.985 and validation perplexity 146.24. The
final checkpoint at step 8500 remained close, with validation loss 4.996 and
validation perplexity 147.87, indicating only mild late-stage degradation.

RoPE also substantially improves over sinusoidal positional encoding. The
completed GQA + RoPE rerun achieved its best logged validation performance at
step 6500, reaching validation loss 5.012 and validation perplexity 150.28.
Although ALiBi slightly outperforms RoPE in the current experiments, both
methods dramatically improve over the GQA + sinusoidal baseline. Relative
positional bias also improves over sinusoidal, reaching validation perplexity
235.86, but it underperforms both ALiBi and RoPE by a wide margin. This suggests
that the simple learned relative-bias formulation used here is less effective
than the more structured extrapolation-friendly encodings in this setting.

\subsection{Length Extrapolation}

To directly test whether the positional encodings generalize beyond their
training context, I trained models with sequence length 512 and evaluated their
final checkpoints at longer validation contexts. The standard Transformer with
sinusoidal encoding is included as the absolute-position baseline, and the GQA
runs compare RoPE, ALiBi, and learned relative position bias.

\begin{table}[h]
\centering
\begin{tabular}{llrrr}
\toprule
Method & Eval Context & Val Loss & Val PPL & Eval Time (s) \\
\midrule
Sinusoidal & 512 & 5.7039 & 300.04 & 7.95 \\
Sinusoidal & 1024 & 6.2477 & 516.83 & 9.10 \\
Sinusoidal & 2048 & 7.1561 & 1281.94 & 11.16 \\
ALiBi & 512 & 5.0253 & 152.21 & 8.17 \\
ALiBi & 1024 & 5.0033 & 148.90 & 9.91 \\
ALiBi & 2048 & 4.9956 & 147.76 & 13.29 \\
RoPE & 512 & 5.0701 & 159.20 & 8.13 \\
RoPE & 1024 & 5.2424 & 189.13 & 9.39 \\
RoPE & 2048 & 5.6356 & 280.24 & 11.86 \\
Relative & 512 & 5.2876 & 197.86 & 9.29 \\
Relative & 1024 & 5.3719 & 215.27 & 12.04 \\
Relative & 2048 & 5.5349 & 253.37 & 17.90 \\
\bottomrule
\end{tabular}
\caption{Length extrapolation for models trained at context length 512.}
\end{table}

The extrapolation result strongly favors ALiBi. Its validation perplexity
remains stable when the evaluation context is extended beyond the training
length, slightly improving from 152.21 at 512 tokens to 147.76 at 2048 tokens.
RoPE is competitive at the training context, reaching 159.20 perplexity at 512
tokens, but degrades to 189.13 at 1024 tokens and 280.24 at 2048 tokens. Learned
relative bias sits between ALiBi and RoPE at the longest context, reaching
253.37 perplexity at 2048 tokens, but is worse than both at the training length.
All three long-context encodings are substantially better than the standard
sinusoidal baseline, which collapses from 300.04 perplexity at 512 tokens to
1281.94 at 2048 tokens. This is consistent with ALiBi's design: the positional
signal is computed as a distance-dependent attention bias rather than relying on
absolute position values, learned distance tables, or rotary phases that may not
extrapolate as smoothly without additional scaling. The result supports using
GQA + ALiBi as the final Core ML architecture, especially when length
extrapolation is part of the evaluation criterion.

\subsection{Latency Benchmark}

I also benchmarked the selected GQA + ALiBi configuration across sequence
lengths to quantify the cost of longer-context evaluation.

\begin{table}[h]
\centering
\begin{tabular}{lrrrr}
\toprule
Model & Seq Len & Tokens/sec & Avg Latency (ms) & Peak Mem (MB) \\
\midrule
GQA + ALiBi & 512 & 25746.3 & 39.77 & 369.0 \\
GQA + ALiBi & 1024 & 27236.4 & 75.19 & 568.1 \\
GQA + ALiBi & 2048 & 20248.2 & 202.29 & 968.5 \\
\bottomrule
\end{tabular}
\caption{Inference latency benchmark for GQA + ALiBi with batch size 2.}
\end{table}

The benchmark shows the expected growth in latency and memory with sequence
length, but the 2048-token setting remains practical on a single GPU in the
assignment environment. Peak allocated memory stays below 1 GB during the
forward-only benchmark at batch size 2, which makes the selected configuration
usable for long-context evaluation even under constrained compute.

The latency benchmark is reported for the final selected architecture rather
than for every attention variant. Comparative latency baselines for standard,
linear, and sliding-window attention are left as a low-cost extension. This
choice prioritizes validating the selected model's long-context behavior and
keeps the final report focused on the strongest configuration.

\section{Convolution-Attention Hybrids}

Two hybrid designs are implemented: a depthwise convolution branch before each
attention layer, and a gated convolutional feed-forward block. These are tested
with the best attention and positional encoding from the earlier ablations.

\begin{table}[h]
\centering
\begin{tabular}{lrrrr}
\toprule
Block Type & Val PPL & Tokens/sec & Peak Memory MB & Notes \\
\midrule
Standard & 146.24 & 28440.73 & 4153.57 & Best GQA + ALiBi checkpoint \\
Conv Before Attention & 173.83 & 20122.53 & 6453.46 & Worse quality and efficiency \\
Gated Conv FFN & 147.50 & 19425.99 & 6832.47 & Competitive quality but much slower \\
\bottomrule
\end{tabular}
\caption{Hybrid architecture results.}
\end{table}

The convolution-augmented variants were evaluated using the strongest attention
and positional configuration from the earlier experiments: GQA with ALiBi. The
standard Transformer block remains the best overall trade-off, achieving the
lowest validation perplexity while also using substantially less memory and
providing higher throughput. Adding a convolution before attention hurts both
quality and efficiency, increasing peak training memory from 4153 MB to 6453 MB
and degrading validation perplexity from 146.24 to 173.83. The gated
convolutional feed-forward block is much more competitive, reaching validation
perplexity 147.50, but it is significantly slower and more memory-intensive than
the standard block. Therefore, the hybrid experiments do not improve the final
model choice, but they show that gated convolution is a plausible local-context
augmentation whereas naive pre-attention convolution is not.

\section{Analysis}

The first major conclusion from the experiments is that not all efficient
attention variants are equally useful in this setting. Sliding-window attention
did not improve either quality or efficiency enough to justify its use at the
tested budget. The implementation uses a full attention matrix with a local
mask, so it should be interpreted as a locality ablation rather than a truly
memory-efficient sparse attention kernel. Linear attention improved validation perplexity slightly over the
standard baseline but required substantially more training memory and reduced
throughput, which weakens its appeal as a practical replacement.

GQA, in contrast, improved both language modeling quality and efficiency. It
reduced validation perplexity from 351.41 to 333.78 while slightly lowering peak
training memory and maintaining very similar throughput. This makes GQA the most
attractive attention choice among the tested variants.

The positional encoding experiments point even more strongly toward the
importance of architecture choices for long-context behavior. ALiBi and RoPE
both produce very large improvements when paired with GQA. ALiBi reduces
validation perplexity from 333.78 for GQA + sinusoidal to a best logged value of
146.24, while RoPE reaches 150.28. Both curves show mild late-stage degradation
after their best checkpoints, suggesting overfitting or diminishing returns
rather than optimization instability. Relative positional bias reaches validation
perplexity 235.86, which is a substantial improvement over sinusoidal but still
well behind ALiBi and RoPE. In the current experiments, ALiBi gives the best
positional encoding result, with RoPE as a close second.

Overall, the experiments suggest a useful hierarchy: GQA is the best attention
mechanism tested so far, and ALiBi is the strongest positional encoding tested so far.
This already gives a strong candidate model for the next stages of the project
and for later hybrid convolution experiments.

The explicit extrapolation experiment strengthens this conclusion. A GQA +
ALiBi model trained at context length 512 remains stable when evaluated at 1024
and 2048 tokens, reaching validation perplexity 147.76 at the longest tested
context. In contrast, a matching GQA + RoPE model trained at 512 tokens degrades
from 159.20 perplexity at 512 tokens to 280.24 at 2048 tokens, learned relative
bias reaches 253.37 at 2048 tokens, and the standard sinusoidal baseline
degrades much more severely from 300.04 to 1281.94. This separates in-context
validation quality from length generalization: ALiBi is not merely the best
positional encoding at the training length, but also the most convincing choice
for extrapolating beyond it.

The convolution-hybrid experiments do not overturn this hierarchy. Gated
convolution nearly matches the standard block's validation perplexity, but its
throughput and memory costs are much worse. The conv-before-attention design is
clearly dominated in the current setup. This suggests that local convolutional
mixing can be useful only when integrated carefully into the feed-forward path,
and even then the marginal quality gain does not justify the compute overhead
for this assignment-scale model.

\section{Conclusion}

The baseline standard Transformer with sinusoidal positional encoding establishes
a reasonable starting point, but it is clearly not the strongest model in the
current study. GQA emerges as the best attention variant, outperforming
standard, sliding-window, and linear attention on validation perplexity while
remaining efficient in both memory and throughput. ALiBi and RoPE further
improve results substantially when combined with GQA, with ALiBi giving the best
observed validation perplexity. The best logged GQA + ALiBi checkpoint occurs at
step 7000, where the model reaches validation perplexity 146.24, while the final
completed run at step 8500 remains close at 147.87. The best GQA + RoPE
checkpoint reaches validation perplexity 150.28, making it a strong second-best
positional encoding result. Convolution-hybrid variants do not improve the best
model: the gated-convolution block is close in perplexity but much slower and
more memory-intensive, while the conv-before-attention design performs worse.

The main lesson from the ablation study is that long-context modeling quality in
this setup depends strongly on choosing the right attention and positional
components rather than only scaling the baseline Transformer. Future work in the
same codebase should add more baseline latency comparisons and, if compute
allows, repeat the 512-to-2048 extrapolation test for RoPE. The current results
already support GQA + ALiBi as the final Core ML submission model.

\section{Reproducibility Artifacts}

The compact experiment artifacts used in this report are versioned under
\texttt{results/core\_ml/}. This includes final metric JSON files, available
training-log tails, extrapolation outputs, latency benchmarks, and an aggregate
\texttt{experiment\_summary.json}. Large checkpoints and raw W\&B/output
directories are intentionally excluded from git to keep the repository small.
The training script also supports mixed precision and gradient accumulation for
future reruns, but the reported experiments use the settings stored in their
logged configs and result artifacts.

\end{document}
